% =============================================================================
% RESUMO
% =============================================================================

\chapter*{Resumo}

As máquinas rotativas desempenham um papel fundamental em diversos setores industriais, sendo sua operação confiável essencial para a eficiência e segurança dos processos. A detecção precoce de falhas, como trincas ou desalinhamentos, é crucial para evitar paradas não planejadas e reduzir custos de manutenção. Este trabalho propõe o desenvolvimento de uma metodologia baseada em simulações numéricas para o diagnóstico de falhas em máquinas rotativas, utilizando a análise de espectros de ordem superior (Higher-Order Spectra - HOS). A metodologia é desenvolvida em três etapas principais: (1) modelagem numérica de sistemas rotativos utilizando a ferramenta ROSS para simular o comportamento dinâmico sob condições normais e com falhas específicas; (2) aplicação de técnicas de HOS aos dados simulados para identificar padrões característicos que diferenciem as condições normais das falhas modeladas; e (3) validação dos resultados obtidos através de comparação com dados disponíveis na literatura. Os resultados demonstram que a análise de bi-espectro e tri-espectro é eficaz para identificar assinaturas não-lineares características de diferentes tipos de falhas em sistemas rotativos. A metodologia proposta oferece uma alternativa baseada em simulação que pode complementar ou reduzir a dependência de experimentos físicos, contribuindo para o aprimoramento das técnicas de diagnóstico de falhas em máquinas rotativas.

\textbf{Palavras-chave:} Máquinas rotativas, Diagnóstico de falhas, Espectros de ordem superior, Simulação numérica, Análise de vibrações.

\cleardoublepage
